\chapter{Objectifs d'apprentissage}\label{II-objectifs-dapprentissage}

\runinhead{Moments d'une variable
aléatoire}\phantomsection\label{II-moments-dune-variable-aluxe9atoire}

\begin{itemize}
\tightlist
\item
  $\mathord{\bullet}$ connaître la définition de l'espérance d'une
  variable aléatoire
\item
  $\mathord{\bullet}$ connaître l'interprétation de l'espérance
\item
  $\mathord{\bullet}$ connaître les propriétés de base de l'espérance
\item
  $\mathord{\bullet}$ connaître la définition de la variance d'une
  variable aléatoire
\item
  $\mathord{\bullet}$ savoir que la variance est toujours positive
\item
  $\mathord{\bullet}$ savoir ce que signifie variable centrée-réduite
\item
  $\mathord{\bullet}$ connaître la définition de la covariance
\item
  $\mathord{\bullet}$ savoir que la covariance est une forme bilinéaire
\item
  $\mathord{\bullet}$ connaître et savoir manipuler l'inégalité de
  Cauchy-Schwarz
\item
  $\mathord{\bullet}\mathord{\bullet}$ connaître les conditions
  d'existence de l'espérance d'une fonction (mesurable) d'une variable
  aléatoire
\item
  $\mathord{\bullet}\mathord{\bullet}$ savoir exprimer la probabilité
  d'un événement de la forme \(X \in A\) sous la forme d'une espérance
\item
  $\mathord{\bullet}\mathord{\bullet}$ connaître et savoir calculer les
  moments des lois usuelles
\end{itemize}

\runinhead{Vecteurs
aléatoires}\phantomsection\label{II-vecteurs-aluxe9atoires}

\begin{itemize}
\tightlist
\item
  $\mathord{\bullet}$ savoir que les composantes d'un vecteur aléatoire
  sont des variables aléatoires définies sur le même espace de
  probabilité
\item
  $\mathord{\bullet}$ connaître la définition d'un vecteur aléatoire à
  densité
\item
  $\mathord{\bullet}\mathord{\bullet}$ connaître les conditions
  d'existence de l'espérance d'une fonction (mesurable) d'un vecteur
  aléatoire à densité
\item
  $\mathord{\bullet}\mathord{\bullet}$ savoir calculer les densités
  marginales d'un vecteur aléatoire à densité
\item
  $\mathord{\bullet}\mathord{\bullet}$ savoir qu'un couple de variables
  aléatoires à densité n'admet pas nécessairement de densité
\item
  $\mathord{\bullet}$ connaître la définition de l'espérance d'un
  vecteur aléatoire à densité
\item
  $\mathord{\bullet}$ connaître la définition de la matrice de
  covariance d'un vecteur aléatoire à densité
\item
  $\mathord{\bullet}$ savoir que la matrice de covariance est
  semi-définie positive
\end{itemize}

\runinhead{Variables et vecteurs aléatoires
indépendantes}\phantomsection\label{II-variables-et-vecteurs-aluxe9atoires-induxe9pendantes}

\begin{itemize}
\tightlist
\item
  $\mathord{\bullet}$ connaître la définition de l'indépendance de deux
  vecteurs/variables aléatoires
\item
  $\mathord{\bullet}$ savoir caractériser l'indépendance de deux
  vecteurs/variables aléatoires à densité
\item
  $\mathord{\bullet}\mathord{\bullet}$ savoir que si \(X\) et \(Y\) sont
  indépendantes, alors \(g(X)\) et \(h(Y)\) sont aussi indépendantes
\item
  $\mathord{\bullet}\mathord{\bullet}$ savoir que si \(X\) et \(Y\) sont
  indépendantes et si \(g(X)\) et \(h(Y)\) sont intégrables alors
  \(\mathbb{E}(g(X)g(Y)) = \mathbb{E}(g(X))\mathbb{E}(g(Y))\)
\item
  $\mathord{\bullet}$ savoir que la covariance de deux variables
  indépendantes est nulle
\item
  $\mathord{\bullet}\mathord{\bullet}$ savoir que la réciproque est
  fausse et connaître un contre-exemple
\end{itemize}

\runinhead{Identification de
densité}\phantomsection\label{II-identification-de-densituxe9}

\begin{itemize}
\tightlist
\item
  $\mathord{\bullet}\mathord{\bullet}$ savoir identifier la densité
  d'une fonction d'une variable aléatoire par la méthode de la fonction
  muette
\item
  $\mathord{\bullet}\mathord{\bullet}$ savoir identifier la densité
  d'une fonction d'un vecteur aléatoire par la méthode de la fonction
  muette
\end{itemize}

\chapter{Moments d'une variable
aléatoire}\label{II-moments-dune-variable-aluxe9atoire-1}

La loi de probabilité d'une variable aléatoire définie sur l'espace
probabilisé \((\Omega,\mathcal{A},\mathbb{P})\) va nous permettre de
calculer aisément des grandeurs caractéristiques telles que sa valeur
moyenne et sa variance (lorsqu'elles existent). Elles sont définies
ci-dessous.

\addcontentsline{toc}{subsection}{Espérance d'une variable aléatoire
réelle}
\begin{definition}[{Espérance d'une variable aléatoire
réelle}]\phantomsection\label{II-defesp}

La variable aléatoire \(X : \Omega \to \mathbb{R}\) de loi
\(\mathbb{P}_X\) est dite \emph{intégrable} si l'intégrale
\(\int_\mathbb{R}|x| d\mathbb{P}_X(x)\) est finie, autrement dit si
\(x\) est \(\mathbb{P}_X\)-intégrable. On définit alors son
\emph{espérance} par
\[\mathbb{E}(X) = \int_\mathbb{R}x d\mathbb{P}_X(x) = \int_\Omega X(\omega) d\mathbb{P}(\omega)\]

Si la variable aléatoire \(X\) admet une densité \(f\) par rapport à la
mesure de Lebesgue, elle sera \emph{intégrable} si l'intégrale
\(\int_\mathbb{R}|x|f(x) dx\) est finie, autrement dit si le produit
\(x f(x)\) est intégrable par rapport à la mesure de Lebesgue. Son
\emph{espérance} vaut alors \[\mathbb{E}(X) = \int_\mathbb{R}x f(x)dx\]
\end{definition}

\addcontentsline{toc}{subsection}{Interprétation}
\begin{remarque}[{Interprétation}]\phantomsection\label{II-interpruxe9tation}

\(\mathbb{E}(X)\) est un nombre réel qui donne une valeur moyenne
résumant la variable aléatoire \(X\).

L'espérance mathématique d'une variable aléatoire est un concept
fondamental de la théorie des probabilités. La dénomination d'espérance
pour cette quantité fait référence aux problèmes de jeux et d'espérance
de gain. Cette terminologie imagée a été introduite par Pascal.

On note \(\mathcal{L}^1\) l'ensemble de toutes les variables réelles
\(X\) intégrables. Les propriétés suivantes découlent directement des
propriétés de l'intégrale.
\end{remarque}

\addcontentsline{toc}{subsection}{Propriétés de base}
\begin{proposition}[{Propriétés de
base}]\phantomsection\label{II-propl1}

\begin{itemize}
\tightlist
\item
  \(\mathcal{L}^1\) est un espace vectoriel et
  \(\forall X,Y \in \mathcal{L}^1\), \(\forall a,b \in \mathbb{R}\)
  \[\mathbb{E}(aX + bY) = a\mathbb{E}(X) + b\mathbb{E}(Y).\]
\item
  \(X \in \mathcal{L}^1 \Leftrightarrow |X| \in \mathcal{L}^1\), et dans
  ce cas \[|\mathbb{E}(X)| \leq \mathbb{E}(|X|).\]
\item
  Si \(X \geq 0\) et \(X \in \mathcal{L}^1\), alors
  \(\mathbb{E}(X) \geq 0.\)
\item
  Si \(X,Y \in \mathcal{L}^1\) sont telles que \(X \leq Y\), alors
  \[ \mathbb{E}(X) \leq \mathbb{E}(Y).\]
\item
  L'espérance d'une variable presque-sûrement constante est égale à
  cette constante~:
  \[ \text{Si } \mathbb{P}(X(\omega) = a ) = 1, \text{ alors } \mathbb{E}(X) = a.\]
\item
  Si \(\exists b \in \mathbb{R}_+\) tel que \(|X| \leq b\), alors
  \(X \in \mathcal{L}^1\) et \(\mathbb{E}(X) \leq b\).
\end{itemize}
\end{proposition}

\runinhead{Exercice -- Démontrer ces
propriétés (\(\mathord{\bullet}\))}\phantomsection\addcontentsline{toc}{subsubsection}{Démontrer
ces propriétés}\label{II-basics}
\addcontentsline{toc}{subsection}{Rappel~: cas discret}

\begin{remarque}[{Rappel~: cas
discret}]\phantomsection\label{II-rappel-cas-discret}
Dans le cas d'une variable aléatoire discrète \(X\) à valeurs dans
\(\mathbb{N}^\ast\), son espérance est définie par la quantité
\(\mathbb{E}(X) = \sum_{i\in\mathbb{N}^\ast} i\mathbb{P}(X=i)\), pourvu
que celle-ci soit finie. On voit immédiatement que les propriétés
ci-dessus sont également vérifiées.
\end{remarque}

Outre l'espace \(\mathcal{L}^1\), nous pouvons définir l'espace
\(\mathcal{L}^2\) des variables aléatoires réelles dont le carré \(X^2\)
est dans \(\mathcal{L}^1\).

\addcontentsline{toc}{subsection}{Variance}
\begin{definition}[{Variance}]\phantomsection\label{II-defvar}

La variable aléatoire \(X : \Omega \to \mathbb{R}\) de loi
\(\mathbb{P}_X\) est dite \emph{de carré intégrable} si
\(\mathbb{E}(X^2) = \int_{\mathbb{R}} x^2 \mathbb{P}_X(dx) = \int_{\Omega} X^2(\omega)\mathbb{P}(d\omega) < +\infty\).
Sa \emph{variance} est définie par
\[\mathbb{V}(X) = \mathbb{E}((X-\mathbb{E}(X))^2)\]

Si la variable aléatoire \(X\) admet une densité \(f\) par rapport à la
mesure de Lebesgue, elle sera \emph{de carré intégrable} si
\(\mathbb{E}(X^2) = \int_\mathbb{R}x^2 f(x)dx\) est définie, autrement
dit si le produit \(x^2 f(x)\) est intégrable par rapport à la mesure de
Lebesgue.
\end{definition}

\addcontentsline{toc}{subsection}{Propriétés}
\begin{proposition}[{Propriétés}]\phantomsection\label{II-propl2}

\(\mathcal{L}^2\) est un sous-espace vectoriel de \(\mathcal{L}^1\), et
si \(X \in \mathcal{L}^2\),
\[|\mathbb{E}(X)| \leq \mathbb{E}(|X|) \leq \sqrt{\mathbb{E}(X^2)}\]
\end{proposition}

\begin{preuve}
Soient \(X\) et \(Y\) deux variables aléatoires réelles de
\(\mathcal{L}^2\) et \(a \in \mathbb{R}\). Comme
\((aX+ Y)^2 \leq 2 a^2 X^2 + 2 Y^2\), alors
\(aX + Y \in \mathcal{L}^2\). Ainsi, \(\mathcal{L}^2\) est un espace
vectoriel.

L'inclusion \(\mathcal{L}^2 \subset \mathcal{L}^1\) découle de
\(|X| \leq 1 + X^2\) et de la \protect\hyperref[II-propl1]{proposition
précédente (p.~\pageref*{II-propl1})} (linéarité).

La première inégalité a déjà été vue
\protect\hyperref[II-propl1]{ci-dessus (p.~\pageref*{II-propl1})}. Pour
la seconde, nous pouvons nous limiter au cas où \(X\) est positive. Soit
alors \(a = \mathbb{E}(X)\) et \(Y = X-a\). Par linéarité, on a
\[ \mathbb{E}(Y^2) = \mathbb{E}(X^2) - 2a \mathbb{E}(X) + a^2 = \mathbb{E}(X^2)-a^2.\]
Et \(\mathbb{E}(Y^2) \geq 0\) par le troisième point de la
\protect\hyperref[II-propl1]{proposition
ci-dessus (p.~\pageref*{II-propl1})}. Par conséquent,
\(\mathbb{E}(X)^2 = a^2 \leq \mathbb{E}(X^2)\) ce qui est le résultat
recherché.\hfill$\blacksquare$
\end{preuve}

\runinhead{Exercice -- Approximation par une
constante (\(\mathord{\bullet}\))}\phantomsection\addcontentsline{toc}{subsubsection}{Approximation
par une constante}\label{II-exo-approx}
Soit \(X\) une v.a. de carré intégrable. Montrer que
\(\mathbb{E}((X - a)^2 )\) atteint son minimum lorsque
\(a = \mathbb{E}(X)\).\solref{II-answer-exo-approx}

\addcontentsline{toc}{subsection}{La variance est positive}
\begin{remarque}[{La variance est
positive}]\phantomsection\label{II-la-variance-est-positive}

En vertu de cette
\protect\hyperref[II-propl2]{proposition (p.~\pageref*{II-propl2})},
\(\mathbb{V}(X)\) est \textbf{positive} et sa racine carrée \(\sigma_X\)
s'appelle l'\emph{écart-type} de \(X\). L'écart-type est une grandeur
qui mesure la moyenne (en un certain sens) de l'écart des valeurs de
\(X\) à sa moyenne, d'où son nom.

On a également \[\mathbb{V}(X) = \mathbb{E}(X^2)-\mathbb{E}(X)^2\] que
l'on obtient en développant \((X-\mathbb{E}(X))^2\). Cette manipulation
anodine est fort utile dans la pratique. On retiendra que ``La variance
est égale à la moyenne des carrés moins le carré de la moyenne''. On
désigne le terme \(\mathbb{E}(X^2)\) par l'expression \emph{moment
d'ordre deux} tandis que la variance est parfois appelée \emph{moment
centré d'ordre deux}.
\end{remarque}

\addcontentsline{toc}{subsection}{Variable centrée réduite}
\begin{remarque}[{Variable centrée
réduite}]\phantomsection\label{II-variable-centruxe9e-ruxe9duite}

D'après ce qui précède, si \(X\) est une variable aléatoire de carré
intégrable, d'espérance \(\mathbb{E}(X)\) et d'écart-type
\(\sigma_X >0\), alors la variable aléatoire
\[\frac{X-\mathbb{E}(X)}{\sigma_X}\] est d'espérance nulle et de
variance 1. On dira qu'une telle variable aléatoire est \emph{centrée et
réduite}.

On peut remarquer que si \(X\) et \(Y\) sont dans \(\mathcal{L}^2\), la
variable aléatoire \(XY\) est dans \(\mathcal{L}^1\), puisque
\(|XY| \leq \frac{1}{2}(X^2+Y^2)\). On peut ainsi définir la
\emph{covariance} de deux variables aléatoires~:
\end{remarque}

\addcontentsline{toc}{subsection}{Covariance}
\begin{definition}[{Covariance}]\phantomsection\label{II-defcov}

Si \(X\) et \(Y\) sont dans \(\mathcal{L}^2\), la variable aléatoire
\((X-\mathbb{E}(X))(Y-\mathbb{E}(Y))\) est intégrable. On appelle la
\emph{covariance} de \(X\) et \(Y\) l'espérance de cette variable
aléatoire et on la note~:
\[\text{Cov}(X,Y) = \mathbb{E}((X-\mathbb{E}(X))(Y-\mathbb{E}(Y))).\] Le
\emph{coefficient de corrélation} des variables aléatoires \(X\) et
\(Y\) est le nombre
\[\rho(X,Y) = \frac{\text{Cov}(X,Y)}{\sqrt{\mathbb{V}(X)}\sqrt{\mathbb{V}(Y)}}\]
qui est bien défini lorsque \(\mathbb{V}(X)>0\) et \(\mathbb{V}(Y)>0\).
\end{definition}

\runinhead{Exercice -- Calcul d'une
covariance (\(\mathord{\bullet}\))}\phantomsection\addcontentsline{toc}{subsubsection}{Calcul
d'une covariance}\label{II-cov-unif}
Soit \(U \sim \mathcal{U}_{[0,1]}\) et \(V = 1-U\). Calculer
\(\text{Cov}(U,V)\) et \(\rho(U,V)\).\solref{II-answer-cov-unif}

Du fait de la linéarité de l'espérance, on a
\[\text{Cov}(X,Y) = \mathbb{E}(XY) - \mathbb{E}(X)\mathbb{E}(Y)\] et
d'ailleurs, on voit que la formule de calcul de la variance donnée plus
haut est un cas particulier de cette formulation car
\(\mathbb{V}(X) = \text{Cov}(X,X)\). La linéarité de l'espérance nous
donne encore pour \(a,a',b,b' \in \mathbb{R}\) \begin{align*} 
\mathbb{E}((aX+b)(a'Y+b')) &= aa'\mathbb{E}(XY) + ab'\mathbb{E}(X) +a'b\mathbb{E}(Y) + bb'\\
\mathbb{E}(aX+b)\mathbb{E}(a'Y+b') &= aa'\mathbb{E}(X)\mathbb{E}(Y) + ab'\mathbb{E}(X) +a'b\mathbb{E}(Y) + bb'
\end{align*} On en déduit que la covariance est une forme bilinéaire sur
l'espace vectoriel des variables aléatoires de carré intégrable, et nous
avons \[\text{Cov}(aX+b,a'Y+b') = aa'\text{Cov}(X,Y)\] En particulier,
on a \[ \mathbb{V}(aX) = a^2\mathbb{V}(X)\] On en déduit que les
coefficients de corrélation de \(X\) et \(Y\) et de \(aX+b\) et
\(a'Y+b'\) sont égaux lorsque \(aa' >0\).

\addcontentsline{toc}{subsection}{Inégalité de Cauchy-Schwarz}
\begin{proposition}[{Inégalité de
Cauchy-Schwarz}]\phantomsection\label{II-CS}

Soit \(X\) et \(Y\) deux variables aléatoires de carré intégrable, alors
on a \emph{l'inégalité de Cauchy-Schwarz}~:
\[|\mathbb{E}(XY)| \leq \mathbb{E}(|XY|) \leq \sqrt{\mathbb{E}(X^2)\mathbb{E}(Y^2)}\]
avec égalité si et seulement si \(X\) et \(Y\) sont presque-sûrement
proportionnelles.
\end{proposition}

\begin{preuve}
La première inégalité a été démontrée plus haut. Pour la seconde, on a
\(\forall x \in \mathbb{R}\) d'après la linéarité de l'espérance~:
\[x^2\mathbb{E}(X^2) + 2x\mathbb{E}(XY)+ \mathbb{E}(Y^2) = \mathbb{E}((xX+Y)^2) \geq 0.\]
Mais ceci n'est possible que si ce trinôme en \(x\) n'a au plus qu'une
seul racine réelle\,; son discriminant
\[4\left((\mathbb{E}(XY))^2 - \mathbb{E}(X^2)\mathbb{E}(Y^2)\right)\]
doit donc être négatif ou nul ce qui donne le résultat.

Le discriminant est nul si et seulement si le trinôme admet une racine
double \(x_0\) et dans ce cas, \(Y(\omega) = -x_0 X(\omega)\) pour
presque tout \(\omega\).\hfill$\blacksquare$
\end{preuve}

\runinhead{Exercice -- Coefficient de
corrélation (\(\mathord{\bullet}\))}\phantomsection\addcontentsline{toc}{subsubsection}{Coefficient
de corrélation}\label{II-CS-corr}
Montrer que le coefficient de corrélation de \(X\) et \(Y\) vérifie
\[-1\leq \rho(X,Y) \leq 1.\]

\solref{II-answer-CS-corr}

Enfin, il peut être intéressant de pouvoir calculer l'espérance d'une
fonction mesurable d'une variable aléatoire réelle qui est une variable
aléatoire en vertu de la
\href{Probabilité\%20I.pdf\%20\#composition}{proposition de
composition}.

\addcontentsline{toc}{subsection}{Espérance de \(g(X)\)}
\begin{proposition}[{Espérance de
\(g(X)\)}]\phantomsection\label{II-esperanceg}

Soit \(X\) une variable aléatoire réelle de loi \(\mathbb{P}_X\), et
\(g\) une fonction mesurable de \((\mathbb{R},\mathcal{B}(\mathbb{R}))\)
dans \((\mathbb{R},\mathcal{B}(\mathbb{R}))\). Alors \(g(X)\) est
\(\mathbb{P}_X\) - intégrable si et seulement si l'intégrale
\[\int_\mathbb{R}|g(x)| d\mathbb{P}_X(x),\] est finie et dans ce cas
\[\mathbb{E}(g(X)) = \int_\mathbb{R}g(x)d\mathbb{P}_X(x) = \int_\Omega g(X(\omega)) d\mathbb{P}(\omega)\]

Si \(\mathbb{P}_X\) admet une densité \(f\) par rapport à la mesure de
Lebesgue, alors \(g(X)\) est \(\mathbb{P}_X\) - intégrable si et
seulement si le produit \(fg\) est intégrable par rapport à la mesure de
Lebesgue. On a alors~:
\[\mathbb{E}(g(X)) =  \int_\mathbb{R}g(x)d\mathbb{P}_X(x) = \int_\mathbb{R}g(x)f(x)dx = \int_\Omega g(X(\omega)) d\mathbb{P}(\omega)\]
\end{proposition}

\begin{preuve}
On rappelle que la loi de probabilité \(\mathbb{P}_X\) est définie pour
\(B\in \mathcal{B}(\mathbb{R})\) par
\(\mathbb{P}_X(B)= \mathbb{P}(X^{-1}(B))\). Par conséquent
\[\mathbb{E}(1_B(X)) = \mathbb{P}(X^{-1}(B)) = \mathbb{P}_X(B) = \int_\mathbb{R}1_B (X) d\mathbb{P}(x).\]
Ainsi, si \(g\) est une fonction étagée positive
(\(g \in \mathcal{E}^+\)), on obtient le résultat voulu par linéarité.

Si \(g\) est positive, soit \((g_n)_{n \in \mathbb{N}^\star}\) une suite
croissante de fonctions étagées positives convergeant simplement vers
\(g\). On a alors \begin{align*}
\mathbb{E}(g(x)) & = \mathbb{E}(\lim_{n\to \infty}g_n(X))\\
           & = \lim_{n\to \infty} \mathbb{E}(g_n(X)) ~\text{ par le théorème de convergence monotone}\\
           & = \lim_{n\to \infty} \int_\mathbb{R}g_n(x) d\mathbb{P}_X(x)\\
           & = \int_\mathbb{R}\lim_{n\to \infty} g_n(x) d\mathbb{P}_X(x) ~\text{ par le théorème de convergence monotone}\\
           & = \int_\mathbb{R}g(x) d\mathbb{P}_X(x).
\end{align*}

Enfin, si \(g\) n'est pas positive, on utilise la décomposition
\(g = g^+ - g^-\) et on déduit le résultat par
soustraction.\hfill$\blacksquare$
\end{preuve}

\addcontentsline{toc}{subsection}{Cas particulier}
\begin{remarque}[{Cas
particulier}]\phantomsection\label{II-cas-particulier}

L'espérance et la variance sont des cas particuliers de ce résultat.
\end{remarque}

\addcontentsline{toc}{subsection}{Inégalité de Jensen}
\begin{proposition}[{Inégalité de
Jensen}]\phantomsection\label{II-jensen}

Soit \(X\) une variable aléatoire réelle intégrable de loi
\(\mathbb{P}_X\), et soit \(g : \mathbb{R}\to \mathbb{R}\) une fonction
convexe telle que \(g(X)\) soit \(\mathbb{P}_X\) - intégrable. On a
alors \[ g(\mathbb{E}(X)) \leq \mathbb{E}(g(X)).\]
\end{proposition}

\begin{preuve}
Puisque \(g\) est convexe, pour tout \(a \in \mathbb{R}\) il existe
\(\lambda_a \in \mathbb{R}\) tel que pour tout \(x\in\mathbb{R}\) on a
\[g(x) \geq g(a) + \lambda_a\,(x-a).\] C'est une conséquence directe de
la caractérisation de la convexité par les inégalités des pentes. C'est
vrai en particulier pour \(x =X(\omega)\), \(\omega \in \Omega\), et
\(a = \mathbb{E}(X)\)~: pour tout \(\omega \in \Omega\),
\[g\left(X(\omega)\right) \geq g\left(\mathbb{E}(X)\right) + \lambda_{\mathbb{E}(X)}\,\left(X(\omega) - \mathbb{E}(X)\right).\]
En intégrant de chaque côté de l'inégalité, on obtient bien
\begin{align*}
\int_\Omega g\left(X(\omega)\right) \,\mathbb{P}(d\omega) = \mathbb{E}\left(g(X)\right) &\geq g\left(\mathbb{E}(X)\right) + \lambda_{\mathbb{E}(X)}\,\left(\int_\Omega X(\omega)\,\mathbb{P}(d\omega) - \mathbb{E}(X)\right)\\
& = g\left(\mathbb{E}(X)\right) + \lambda_{\mathbb{E}(X)}\,\left(\mathbb{E}(X) - \mathbb{E}(X)\right)\\
& = g\left(\mathbb{E}(X)\right).
\end{align*}\hfill$\blacksquare$
\end{preuve}

\section{Exemples de loi à
densité}\label{II-exemples-de-loi-uxe0-densituxe9}

Nous donnons ici quelques exemples de densités de probabilité. Nous
reprenons en particulier les
\href{Probabilité\%20I.pdf\%20\#exampledens}{trois exemples de densités
donnés au premier cours}~:

\subsection{\emph{Loi uniforme}}\label{II-loi-uniforme}

sur \([a,b]\), où \(a < b\) et on note \(X \sim \mathcal{U}_{[a,b]}\) si
\(X\) est de densité \[ f(x) = \frac{1}{b-a} 1_{[a,b]} (x).\]

\runinhead{Exercice -- Moments d'une variable aléatoire de loi uniforme
sur
\([a,b]\) (\(\mathord{\bullet}\))}\phantomsection\addcontentsline{toc}{subsubsection}{Moments
d'une variable aléatoire de loi uniforme sur
\([a,b]\)}\label{II-moments-unif}
Soit \(X \sim \mathcal{U}_{[a,b]}\). Montrer que son espérance vaut
\[ \mathbb{E}(X) = \int_a^b \frac{x}{b-a} dx = \frac{a+b}{2}\] et que sa
variance vaut
\[ \mathbb{V}(X) = \mathbb{E}(X^2) - \mathbb{E}(X)^2 = \frac{(b-a)^2}{12}.\]

\solref{II-answer-moments-unif}

\subsection{\emph{Loi exponentielle}}\label{II-loi-exponentielle}

de paramètre \(\theta > 0\) et on note \(X \sim \mathcal{E}(\theta)\) si
\(X\) est de densité \[ f(x) = \theta e^{-\theta x} 1_{\{x>0\}} (x).\]

\runinhead{Exercice -- Moments d'une v.a.
exponentielle (\(\mathord{\bullet}\))}\phantomsection\addcontentsline{toc}{subsubsection}{Moments
d'une v.a. exponentielle}\label{II-moments-expo}
Montrer que son espérance et sa variance valent
\[ \mathbb{E}(X) = \frac{1}{\theta} \text{ et } \mathbb{V}(X) = \frac{1}{\theta^2}\]

\solref{II-answer-moments-expo}

\subsection{\emph{Loi gamma}}\label{II-loi-gamma}

On rappelle tout d'abord que la fonction gamma est définie pour
\(\alpha \in \left] 0, + \infty \right[\) par
\[\Gamma(\alpha) = \int_0^{+\infty} x^{\alpha-1}e^{-x} dx.\] En
intégrant par partie, on obtient la relation
\(\Gamma(\alpha+1) = \alpha \Gamma(\alpha)\), et on a \(\Gamma(1) = 1\).
On en déduit que \(\Gamma(n) = (n-1)!\).

Une variable aléatoire \(X\) suit une loi gamma de paramètre d'échelle
\(\theta\) et d'indice \(\alpha\) et on note
\(X \sim \Gamma(\alpha,\theta)\), si sa loi admet la densité
\[f(x) = \frac{1}{\Gamma(\alpha)}\theta^\alpha x^{\alpha -1} e^{-\theta x}.\]

\runinhead{Exercice -- Moments d'une v.a. de loi
gamma (\(\mathord{\bullet}\))}\phantomsection\addcontentsline{toc}{subsubsection}{Moments
d'une v.a. de loi gamma}\label{II-moments-gamma}
Soit \(X\sim \Gamma(\alpha,\theta)\). Montrer que son espérance et sa
variance valent~:

\[ \mathbb{E}(X) = \frac{\alpha}{\theta} \text{ et } \mathbb{V}(X) = \frac{\alpha}{\theta^2}.\]

On remarquera que \(\Gamma(1,\theta)\) est la loi exponentielle de
paramètre \(\theta\).

Lorsque \(\alpha\) est entier, la loi gamma permet de modéliser le temps
d'attente avant la \(n\)-ième occurence d'événements indépendants de loi
exponentielle de paramètre \(\theta\).\solref{II-answer-moments-gamma}

\subsection{\emph{Loi normale}}\label{II-loi-normale}

de paramètres \(\mu\) et \(\sigma^2\) et on note
\(X \sim \mathcal{N}(\mu,\sigma^2)\) si \(X\) est de densité
\[f(x) = \frac{1}{\sqrt{2\pi}\sigma}\exp\left(-\frac{(x-\mu)^2}{2\sigma^2}\right)\]
Son espérance et sa variance valent
\[\mathbb{E}(X) =  \mu \text{ et } \mathbb{V}(X) = \sigma^2\]

\runinhead{Exercice -- Moments d'une v.a.
gaussienne (\(\mathord{\bullet}\))}\phantomsection\addcontentsline{toc}{subsubsection}{Moments
d'une v.a. gaussienne}\label{II-moments-gauss}
Après avoir vérifié que \(f\) est bien une densité, calculer l'espérance
et la variance de
\(X \sim \mathcal{N}(\mu,\sigma^2)\)\solref{II-answer-moments-gauss}

\addcontentsline{toc}{subsection}{Estimation statistique}
\begin{remarque}[{Estimation
statistique}]\phantomsection\label{II-estimation-statistique}

Dans les exemples ci-dessus, on peut remarquer que les densités sont
paramétrées par des nombres réels qui sont liés directement aux valeurs
de l'espérance et de la variance de la variable. C'est très utile en
statistique où l'on cherchera à estimer la valeur de ces paramètres à
partir des observations disponibles. Dans le cas de la loi normale, la
moyenne et la variance (empiriques) des échantillons fourniront ainsi
directement des estimateurs des paramètres.

Il existe des variables aléatoires qui n'ont pas d'espérance, comme le
montre l'exercice suivant.
\end{remarque}

\runinhead{Exercice -- \emph{Loi de
Cauchy} (\(\mathord{\bullet}\mathord{\bullet}\))}\phantomsection\addcontentsline{toc}{subsubsection}{\emph{Loi
de Cauchy}}\label{II-exo-cauchy}
Un gyrophare envoie un flash lumineux dans une direction aléatoire
uniforme d'angle \(\theta\). On cherche la distribution de l'abscisse
\(X\) du point d'impact du rayon lumineux sur un écran plan infini situé
à distance 1 du gyrophare. Donner la densité de \(X\). Calculer son
espérance.\solref{II-answer-exo-cauchy}

\chapter{Rappels sur les intégrales
multiples}\label{II-rappels-sur-les-intuxe9grales-multiples}

Dans la suite, nous aurons besoin de quelques résultats à propos des
intégrales multiples que nous reformulons ici dans le contexte des
probabilités.

\addcontentsline{toc}{subsection}{Théorème de Fubini (pour les
probabilités)}
\begin{theoreme}[{Théorème de Fubini (pour les
probabilités)}]\phantomsection\label{II-fubiniproba}

Soient \(\mathbb{P}_X\) et \(\mathbb{P}_Y\) deux probabilités définies
respectivement sur \((\mathbb{R}^n, \mathcal{B}(\mathbb{R}^n))\) et
\((\mathbb{R}^m, \mathcal{B}(\mathbb{R}^m))\).

\begin{enumerate}
\def\labelenumi{\arabic{enumi}.}
\tightlist
\item
  Soient \(A \in \mathcal{B}(\mathbb{R}^n)\) et
  \(B \in \mathcal{B}(\mathbb{R}^m)\), alors
  \[\mathbb{P}(A \times B) = \mathbb{P}_X(A)\mathbb{P}_Y(B)\] définit de
  manière unique une probabilité sur
  \((\mathbb{R}^{n+m}, \mathcal{B}(\mathbb{R}^{n+m}))\) (que l'on note
  aussi \(\mathbb{P}_X \otimes \mathbb{P}_Y\)).
\item
  Soit \(f\) une fonction \(\mathcal{B}(\mathbb{R}^{n+m})\)-mesurable
  positive ou \(\mathbb{P}_X \otimes \mathbb{P}_Y\)-intégrable, alors la
  fonction \(x \mapsto \int f(x, y)\mathbb{P}_Y(dy)\) est
  \(\mathcal{B}(\mathbb{R}^{n})\)-mesurable, la fonction
  \(y \mapsto \int f (x, y) \mathbb{P}_X(dx)\) est
  \(\mathcal{B}(\mathbb{R}^m)\)-mesurable et
  \[\int f d\mathbb{P}_X\otimes\mathbb{P}_Y = \int \left(\int f(x,y) \mathbb{P}_Y(dy) \right) \mathbb{P}_X(dx) = \int \left(\int f(x,y) \mathbb{P}_X(dx) \right) \mathbb{P}_Y(dy)\]
\end{enumerate}
\end{theoreme}

\runinhead{Exercice -- Calcul de l'aire d'un
triangle (\(\mathord{\bullet}\))}\phantomsection\addcontentsline{toc}{subsubsection}{Calcul
de l'aire d'un triangle}\label{II-triangle}
Considérons le triangle \[
T = \{(x, y) \in \mathbb{R}^2 \; | \; x \geq 0, \, y \geq 0 \mbox{ et } x + y \leq  1\}.
\] En supposant l'intégrale ci-dessous bien définie, calculer~: \[
a(T) = \int_{\mathbb{R}^2} 1_T(x, y) \, dxdy.
\]

\solref{II-answer-triangle}

\runinhead{Exercice -- Contre-exemple (\(\mathord{\bullet}\mathord{\bullet}\))}\phantomsection\addcontentsline{toc}{subsubsection}{Contre-exemple}\label{II-fubini-counter-example}
Comparer les valeurs des intégrales multiples \[
\int_0^1 \left[\int_0^1 \frac{x^2 - y^2}{(x^2+y^2)^2} \, dx \right] \, dy
\; \mbox{ et } \;
\int_0^1 \left[\int_0^1 \frac{x^2 - y^2}{(x^2+y^2)^2} \, dy \right] \, dx,
\] puis expliquer le résultat. Indication~: on remarquera que \[
\frac{x^2 - y^2}{(x^2+y^2)^2} = \frac{\partial}{\partial x} \left(-\frac{x}{x^2+y^2}\right) 
\] \solref{II-answer-fubini-counter-example}

Pour pouvoir appliquer le théorème de Fubini, il faut savoir a priori
que la fonction est intégrable et pas seulement que ses intégrales
multiples sont bien définies
(cf.~\protect\hyperref[II-fubini-counter-example]{``Contre-exemple''
ci-dessus (p.~\pageref*{II-fubini-counter-example})}). Si la fonction
est à valeurs positives toutefois, l'examen de ses intégrales multiples
permet de s'assurer de l'intégrabilité\,; c'est le théorème de
Fubini-Tonelli vu dans le cours de calcul intégral.

Autrement dit, en pratique, on appliquera Fubini-Tonelli à la valeur
absolue de la fonction pour vérifier que celle-ci est intégrable, puis
Fubini, s'il est applicable, pour effectuer le calcul de l'intégrale.

\addcontentsline{toc}{subsection}{Changement de variables}
\begin{theoreme}[{Changement de
variables}]\phantomsection\label{II-theorem-changement-de-variables}

Soient \(D_1\) et \(D_2\) des ouverts de \(\mathbb{R}^n\) et
\(h: D_1 \to D_2\) un \(C^1\)-difféomorphisme de \(D_1\) sur \(D_2\),
c'est-à-dire une fonction continûment différentiable et bijective dont
l'inverse \(h^{-1}: D_2 \to D_1\) est également continûment
différentiable. La matrice jacobienne associée à la différentielle de
\(h\) étant notée \(J_h\), la fonction \(f: D_2 \to \mathbb{R}\) est
intégrable si et seulement si la fonction
\((f \circ h) |\det J_h| : D_1 \to \mathbb{R}\) est intégrable et dans
ce cas, \[
\int_{D_2} f(y) \, dy = \int_{D_1} f(h(x)) |\det J_h(x)| \, dx.
\]
\end{theoreme}

\begin{figure}
\centering
\fbox{[figure manquante : changement-de-variables.svg.pdf]}
\caption{Changement de
variables}\label{II-fig-changement-de-variables.svg}
\end{figure}

\runinhead{Exercice -- Homothétie (\(\mathord{\bullet}\))}\phantomsection\addcontentsline{toc}{subsubsection}{Homothétie}\label{II-h}
Soit \(f:\mathbb{R}^n \to \mathbb{R}\) une fonction intégrable. Montrer
que pour tout coefficient \(\alpha > 0\), l'intégrale \[
\int_{\mathbb{R}^n} f(\alpha x) dx
\] est bien définie et la calculer en fonction de l'intégrale de \(f\)
sur \(\mathbb{R}^n\).\solref{II-answer-h}

\runinhead{Exercice -- Volume et
translation (\(\mathord{\bullet}\))}\phantomsection\addcontentsline{toc}{subsubsection}{Volume
et translation}\label{II-vr}
Montrer que si \(A\) est mesurable et de volume fini dans
\(\mathbb{R}^3\) l'image de \(A\) par une translation est également
mesurable et de même volume.\solref{II-answer-vr}

\runinhead{Exercice -- Coordonnées
polaires (\(\mathord{\bullet}\mathord{\bullet}\))}\phantomsection\addcontentsline{toc}{subsubsection}{Coordonnées
polaires}\label{II-cp}
Soit \[
C = \{(x, y) \in \mathbb{R}^2 \; | \; y \neq 0 \mbox{ ou } x > 0\}
\; \mbox{ et } \;
P = \{(r,\theta) \in \mathbb{R}^2 \; | \; r>0 \mbox{ et } -\pi < \theta < \pi\}.\]
On note \(h\) la fonction de \(P\) dans \(C\) définie par
\(h(r, \theta) = (r \cos \theta, r \sin \theta)\). Montrer que pour
toute fonction \(f: C \to \mathbb{R}\) intégrable, si l'on pose
\(g(r,\theta) = f(x, y)\) où \((x, y) = h(r,\theta)\), alors \[
\int_{C} f(x, y) \, d(x,y) = \int_{P} g(r,\theta)  r \, d(r, \theta).
\]

\solref{II-answer-cp}

\runinhead{Exercice -- Absence du déterminant
jacobien (\(\mathord{\bullet}\mathord{\bullet}\))}\phantomsection\addcontentsline{toc}{subsubsection}{Absence
du déterminant jacobien}\label{II-adj}
Supposons \(D_1\), \(D_2\), \(h\) et \(f\) conformes aux hypothèses
\protect\hyperref[II-theorem-changement-de-variables]{du théorème de
changement de
variables (p.~\pageref*{II-theorem-changement-de-variables})}. On
suppose de plus que \(f \circ h\) est intégrable sur \(D_1\). Exprimer
l'intégrale \[
\int_{D_1} f (h(x)) \, dx
\] comme une intégrale sur \(D_2\).\solref{II-answer-adj}

\chapter{Vecteurs aléatoires}\label{II-vecteurs-aluxe9atoires-1}

Nous allons généraliser ici la notion de variable aléatoire en
considérant qu'elle peut prendre ses valeurs dans \(\mathbb{R}^n\). Les
vecteurs aléatoires se rencontrent naturellement lorsqu'on s'intéresse à
plusieurs quantités conjointement, par exemple dans le cas de la météo,
la température, la pluviométrie et la vitesse et la direction du vent.

\section{Définitions}\label{II-duxe9finitions}

Une variable aléatoire \(X\) à valeurs dans \(\mathbb{R}^n\) (ou vecteur
aléatoire) est simplement une collection de \(n\) variables réelles
définies sur le même espace probabilisé
\((\Omega, \mathcal{A}, \mathbb{P})\), qui sont les \emph{composantes}
de \(X\). On écrit \(X = (X_1,\ldots,X_n)\).

De même qu'en dimension 1, la loi de \(X\) est caractérisée par la
fonction de répartition multi-dimensionnelle
\(F : \mathbb{R}^n \to \mathbb{R}\) définie par
\[F(x_1,\ldots,x_n) = \mathbb{P}_X(X_1\leq x_1,\ldots,X_n\leq x_n)\]
Cependant, caractériser les fonctions de répartition sur
\(\mathbb{R}^n\) est délicat, de sorte que cette notion est rarement
utilisée. Nous allons plus particulièrement nous intéresser aux vecteurs
aléatoires à densité. On gardera tout de même à l'esprit que les lois de
probabilités définies sur \((\mathbb{R}^n,\mathcal{B}(\mathbb{R}^n))\)
n'admettent pas toutes une densité\,!

\addcontentsline{toc}{subsection}{Vecteur aléatoire à densité}
\begin{definition}[{Vecteur aléatoire à
densité}]\phantomsection\label{II-defvect}

On dit que \(X\) admet la densité \(f\) si la fonction réelle \(f\) sur
\(\mathbb{R}^n\) est positive, intégrable et vérifie
\[\int_{\mathbb{R}^n} f(x) dx = \int_{-\infty}^{+\infty} \ldots \int_{-\infty}^{+\infty} f(x_1,\ldots,x_n) dx_1 \ldots dx_n= 1\]
et si
\[\mathbb{P}_X(X_1\leq x_1,\ldots,X_n\leq x_n) = \int_{-\infty}^{x_1} \ldots \int_{-\infty}^{x_n} f(x_1,\ldots,x_n) dx_1 \ldots dx_n.\]

De la même manière que dans le \protect\hyperref[II-esperanceg]{cas
unidimensionnel (p.~\pageref*{II-esperanceg})}, on a~:
\end{definition}

\addcontentsline{toc}{subsection}{Fonction d'un vecteur aléatoire}
\begin{proposition}[{Fonction d'un vecteur
aléatoire}]\phantomsection\label{II-esperancegvect}

Soit \(X\) un vecteur aléatoire de densité \(f\), et soit \(g\) une
fonction de \(\mathbb{R}^n\) dans \(\mathbb{R}\), mesurable. On a alors
\(g(X)\in \mathcal{L}^1\) si et seulement si
\[\int_{-\infty}^{+\infty} \ldots \int_{-\infty}^{+\infty} |g(x_1,\ldots,x_n)|f(x_1,\ldots,x_n) dx_1 \ldots dx_n,\]
est finie et dans ce cas, on a
\[\mathbb{E}(g(X)) = \int_{-\infty}^{+\infty} \ldots \int_{-\infty}^{+\infty} g(x_1,\ldots,x_n)f(x_1,\ldots,x_n) dx_1 \ldots dx_n.\]

Pour revenir à la densité d'une composante d'un vecteur aléatoire, on
intègre par rapport aux autres variables. On le présente ici dans le cas
d'un couple \(Z=(X,Y)\) de variables aléatoires. On généralise aisément
à une dimension supérieure.
\end{proposition}

\addcontentsline{toc}{subsection}{Densités marginales}
\begin{proposition}[{Densités
marginales}]\phantomsection\label{II-loimarg}

Supposons que \(Z\) admette une densité \(f\). Alors \(X\) et \(Y\)
admettent les densités \(f_X\) et \(f_Y\) données par
\[f_X(x) = \int_\mathbb{R}f(x,y) dy,\,\,\, f_Y(y) = \int_\mathbb{R}f(x,y) dx.\]

Les fonctions \(f_X\) et \(f_Y\) s'appellent les \emph{densités
marginales} de \(f\).
\end{proposition}

\begin{preuve}
Pour tout \(x \in \mathbb{R}\), on a par définition
\[ \mathbb{P}(X \leq x) = \mathbb{P}(Z \in ]-\infty,x] \times \mathbb{R}) = \mathbb{E}(1_{]-\infty,x] \times \mathbb{R}}(Z)) = \int_{-\infty}^x du \left(\int_\mathbb{R}f(u,v) dv\right),\]
où l'on a utilisé le théorème de Fubini-Tonelli pour
\(1_{]-\infty,x] \times \mathbb{R}}(u,v)f(u,v)\) qui est bien positive.
Donc si \(f_X\) est définie par \(f_X(x) = \int_\mathbb{R}f(x,y) dy\),
on obtient que \(\mathbb{P}(X \leq x) = \int_{-\infty}^x f_X(u) du\), ce
qui montre que \(f_X\) est la densité de \(X\). Le raisonnement est
analogue pour \(Y\).\hfill$\blacksquare$
\end{preuve}

\addcontentsline{toc}{subsection}{Réciproque\,?}
\begin{remarque}[{Réciproque\,?}]\phantomsection\label{II-ruxe9ciproque}

La réciproque de cette proposition est fausse en revanche~: les
variables aléatoires \(X\) et \(Y\) peuvent avoir des densités sans que
le couple \(Z = (X, Y )\) en ait une.

Supposons par exemple que \(X = Y\). Si
\(\Delta = \{(x,x) ; x\in \mathbb{R}\}\) est la diagonale de
\(\mathbb{R}^2\), nous avons évidemment \(\mathbb{P}_Z(\Delta) = 1\)
mais si la \protect\hyperref[II-esperancegvect]{proposition
précédente (p.~\pageref*{II-esperancegvect})} était valide pour
\(\mathbb{P}_Z\), on aurait
\(\mathbb{P}_Z(\Delta) = \mathbb{E}(1_\Delta) = \int_{\mathbb{R}^2} 1_\Delta f_Z(z)dz = 0\)
car \(\Delta\) est de volume nul dans \(\mathbb{R}^2\).

En particulier, il faut faire attention au fait que dans le cas général,
la densité d'un couple de variables aléatoires n'est pas le produit des
densités.
\end{remarque}

\runinhead{Exercice -- Lancer de
fléchette (\(\mathord{\bullet}\mathord{\bullet}\))}\phantomsection\addcontentsline{toc}{subsubsection}{Lancer
de fléchette}\label{II-exo-flechette}
On lance une fléchette sur une cible circulaire de rayon unité. Le
joueur est suffisamment maladroit pour que le point \(M\) d'impact de la
fléchette soit supposé uniformément distribué sur la cible (On décide de
n'observer que les lancés qui atteignent la cible). Exprimer la densité
du couple formé par les coordonnées cartésiennes \((X,Y)\) de \(M\) puis
leurs densités marginales.\solref{II-answer-exo-flechette}

\section{Moments d'un vecteur
aléatoire}\label{II-moments-dun-vecteur-aluxe9atoire}

\addcontentsline{toc}{subsection}{Vecteur espérance et covariance}
\begin{definition}[{Vecteur espérance et
covariance}]\phantomsection\label{II-vecteur-espuxe9rance-et-covariance}

Si les composantes \(X_i\) du vecteur aléatoire \(X = (X_1,\ldots,X_n)\)
sont intégrables, on peut définir le \emph{vecteur espérance}
\[\mathbb{E}(X) = (\mathbb{E}(X_1),\ldots,\mathbb{E}(X_n))\] Si les
composantes \(X_i\) du vecteur aléatoire \(X = (X_1,\ldots,X_n)\) sont
de carré intégrable, la \emph{matrice de covariance} de \(X\) est la
matrice \(C_X = (c_{i,j})_{1 \leq i \leq n , 1 \leq j \leq n}\) de
taille \(n \times n\) et dont les éléments valent
\[c_{i,j} = \text{Cov}(X_i,X_j)\]
\end{definition}

\addcontentsline{toc}{subsection}{Matrice de covariance}
\begin{proposition}[{Matrice de
covariance}]\phantomsection\label{II-matrice-de-covariance}

La matrice de covariance est symétrique non-négative (ou encore
semi-définie positive).
\end{proposition}

\begin{preuve}
La symétrie est évidente. Non-négative signifie que pour tous réels
\(a_1,\ldots,a_n\), on a
\(\sum_{i=1}^n \sum_{j=1}^n a_i a_j c_{i,j} \geq 0\). Un calcul simple
montre que
\[ \sum_{i=1}^n \sum_{j=1}^n a_i a_j c_{i,j} = \mathbb{V}(\sum_{i=1}^n a_i X_i) \geq 0.\]

\hfill$\blacksquare$
\end{preuve}

\addcontentsline{toc}{subsection}{Vecteur Gaussien \(n\)-dimensionel}
\begin{exemple}[{Vecteur Gaussien
\(n\)-dimensionel}]\phantomsection\label{II-vecteur-gaussien-n-dimensionel}

Un exemple important de vecteurs aléatoires est celui des vecteurs
gaussiens, que nous étudierons plus en détail au chapitre suivant.
Soient \(m \in \mathbb{R}^n\) et \(C\) une matrice symétrique définie
positive (c'est-à-dire telle que pour tout \(x \in \mathbb{R}^n\) non
identiquement nul \(x^tCx > 0\) où \(^t\) désigne la transposée). Le
vecteur \(X \in \mathbb{R}^n\) est un vecteur aléatoire gaussien
d'espérance \(m\) et de matrice de covariance \(C\) si sa densité
s'écrit
\[ f(x) = \frac{1}{(2\pi)^{n/2}\sqrt{\det (C)}}\exp (-\frac{1}{2}(x-m)^tC^{-1}(x-m)) \]
On a alors \(\mathbb{E}(X) = m\) et \(C_X =C\).
\end{exemple}

\begin{figure}
\centering
\includegraphics{PdfGauss3D.tex.pdf}
\caption{Densité de probabilité de la loi normale bivariée centrée,
réduite, de coefficient de corrélation 1/2}\label{II-fig-PdfGauss3D.tex}
\end{figure}

\chapter{Variables aléatoires
indépendantes}\label{II-variables-aluxe9atoires-induxe9pendantes}

Lorsqu'on modélise plusieurs variables conjointement, une hypothèse
importante est celle de l'indépendance. Ce caractère traduit l'absence
de lien de causalité entre les variables. Par exemple, on fait
naturellement l'hypothèse d'indépendance lorsque l'on considère une
répétition d'une même expérience dans les mêmes conditions.

Dans ce paragraphe, on considère un couple \((X,Y)\) de vecteurs
aléatoires respectivement à valeurs dans \(\mathbb{R}^m\) et
\(\mathbb{R}^n\). Les résultats s'étendent sans peine à une famille
finie quelconque.

On peut se ramener aux évènements pour caractériser l'indépendance de
deux vecteurs aléatoires. En effet, considérons le vecteur aléatoire
\(Z = (X,Y)\), \(A\) et \(B\) deux ensembles dans
\(\mathcal{B}({\mathbb{R}^m})\) et \(\mathcal{B}({\mathbb{R}^n})\). On a
vu que les évènements \(X\in A\) et \(Y \in B\) sont indépendants si et
seulement si
\(\mathbb{P}_Z(X \in A, Y \in B) = \mathbb{P}(X^{-1}(A) \cap Y^{-1}(B)) = \mathbb{P}(X^{-1}(A))\mathbb{P}(Y^{-1}(B)) = \mathbb{P}_X(X \in A)\mathbb{P}_Y(Y \in B)\).
Pour que deux vecteurs aléatoires soient indépendants, on va donc
demander que ceci soit valable quels que soient \(A\) et \(B\).

\addcontentsline{toc}{subsection}{Vecteurs aléatoires indépendants}
\begin{definition}[{Vecteurs aléatoires
indépendants}]\phantomsection\label{II-defvai}

Les vecteurs aléatoires \(X\) et \(Y\) sont \emph{indépendants} si pour
tous ensembles \(A\) et \(B\) des tribus correspondantes,
\[\mathbb{P}(X \in A, Y \in B) = \mathbb{P}(X \in A)\mathbb{P}(Y \in B)\]
\end{definition}

En termes de mesures de probabilité, cette égalité s'écrit
\[\mathbb{P}_{(X,Y)}(A \times B) = \mathbb{P}_X(A)\mathbb{P}_Y(B) = (\mathbb{P}_X \otimes \mathbb{P}_Y) (A \times B).\]
Autrement dit, \(X\) et \(Y\) sont indépendantes si et seulement si la
mesure de probabilité du couple \((X,Y)\) est la mesure produit
\(P_X \otimes P_Y\) dont l'existence et l'unicité sont garanties par le
théorème de
\protect\hyperref[II-fubiniproba]{Fubini (p.~\pageref*{II-fubiniproba})}.

Cette définition se traduit dans le cas à densité dans la proposition
suivante que l'on énonce sans perte de généralité pour un couple de
variables aléatoires.

\addcontentsline{toc}{subsection}{Caractérisation de l'indépendance --
cas à densité}
\begin{proposition}[{Caractérisation de l'indépendance -- cas à
densité}]\phantomsection\label{II-caractuxe9risation-de-linduxe9pendance-cas-uxe0-densituxe9}

Soient \(X\) et \(Y\) deux variables aléatoires réelles de densités
\(f_X\) et \(f_Y\). \(X\) et \(Y\) sont indépendantes si et seulement si
le couple \(Z = (X,Y)\) a pour densité (sur \(\mathbb{R}^2\))~:
\[f_Z(x,y) = f_X(x)f_Y(y).\]
\end{proposition}

\begin{preuve}
S'il y a indépendance, la
\protect\hyperref[II-defvai]{définition (p.~\pageref*{II-defvai})}
implique
\[P(X\leq x, Y\leq y) = P(X\leq x) \mathbb{P}(Y\leq y) = \int_{-\infty}^x f_X(u) du \int_{-\infty}^y f_Y(v) dv\]
ce qui montre que \(\mathbb{P}_Z\) vérifie la
\protect\hyperref[II-defvect]{définition d'un vecteur aléatoire à
densité (p.~\pageref*{II-defvect})} avec \(f_Z=f_X f_Y\).

Inversement, si \(f_Z=f_X f_Y\), on a pour tous \(A\), \(B\) de
\(\mathcal{B}({\mathbb{R}})\)
\[\mathbb{P}(X\in A, Y\in B) = \int_A \int_B f_X(x) f_Y(y) dxdy = \mathbb{P}(X \in A)\mathbb{P}(Y \in B)\]
où on a utilisé le théorème de Fubini-Tonelli.\hfill$\blacksquare$
\end{preuve}

\runinhead{Exercice -- Lancer de fléchette
(suite) (\(\mathord{\bullet}\))}\phantomsection\addcontentsline{toc}{subsubsection}{Lancer
de fléchette (suite)}\label{II-exo-flechettebis}
Les coordonnées cartésiennes \((X,Y)\) de \(M\) sont elles
indépendantes\,?\solref{II-answer-exo-flechettebis}

Considérons maintenant deux fonctions mesurables \(g\) et \(h\) définies
sur \(\mathbb{R}^m\) et \(\mathbb{R}^n\).

\addcontentsline{toc}{subsection}{Fonctions de vecteurs aléatoires
indépendants}
\begin{proposition}[{Fonctions de vecteurs aléatoires
indépendants}]\phantomsection\label{II-indep_fct}

Avec les notations précédentes, si \(X\) et \(Y\) sont indépendantes de
densités respectives \(f_X\) et \(f_Y\), les variables aléatoires
\(g(X)\) et \(h(Y)\) sont aussi indépendantes. Si de plus \(g(X)\) et
\(h(Y)\) sont intégrables, alors le produit \(g(X)h(Y)\) est aussi
intégrable, et on a
\[ \mathbb{E}(g(X)h(Y)) = \mathbb{E}(g(X))\mathbb{E}(h(Y))\]
\end{proposition}

\begin{preuve}
La première assertion est évidente par définition de l'indépendance. Par
ailleurs, si \(g(X)\) et \(h(Y)\) sont intégrables, en notant
\(f_{(X,Y)}\) la densité du couple \((X,Y)\), et en utilisant le
théorème de Fubini, on a \begin{align*}
\mathbb{E}(g(X)h(Y))  & = \int_{\mathbb{R}^{m+n}} g(x)h(y) f_{(X,Y)}(x,y) dx dy \\
                & = \int_{\mathbb{R}^m} \int_{\mathbb{R}^n} g(x)h(y)f_X(x) f_Y(y) dx dy \\
                & = \left(\int_{\mathbb{R}^m}  g(x) f_X(x) dx \right) \left(\int_{\mathbb{R}^n} h(y) f_Y(y) dy\right)\\
                & = \mathbb{E}(g(X))\mathbb{E}(h(Y))
\end{align*}\hfill$\blacksquare$
\end{preuve}

\addcontentsline{toc}{subsection}{Cas général}
\begin{remarque}[{Cas
général}]\phantomsection\label{II-cas-guxe9nuxe9ral}

Ce résultat est encore valable dans le cas général (sans densité). Il
suffit de remplacer \(f_{(X,Y)}(x,y) dx dy\) par
\(d\mathbb{P}_{X,Y}(x,y) = d\mathbb{P}_X(x)d\mathbb{P}_Y(y)\) dans la
démonstration précédente.

On déduit de ce résultat et de la
\protect\hyperref[II-defcov]{définition de la
covariance (p.~\pageref*{II-defcov})} que~:
\end{remarque}

\addcontentsline{toc}{subsection}{Covariance de variables aléatoires
indépendantes}
\begin{corollaire}[{Covariance de variables aléatoires
indépendantes}]\phantomsection\label{II-covariance-de-variables-aluxe9atoires-induxe9pendantes}

Si les variables aléatoires réelles \(X\) et \(Y\) sont indépendantes et
de carré intégrable, alors \(\text{Cov}(X,Y) = 0\) et \(\rho(X,Y) = 0\).
\end{corollaire}

\runinhead{Exercice -- Démonstration (\(\mathord{\bullet}\))}\phantomsection\addcontentsline{toc}{subsubsection}{Démonstration}\label{II-demo-cov-indep}
Démontrer ce résultat.\solref{II-answer-demo-cov-indep}

\subsection{Remarque}\label{II-remarque}

Attention, la réciproque est fausse. Par exemple, si
\(X \sim \mathcal{U}_{[-1,1]}\) et \(Y = X^2\). \(X\) et \(Y\) ne sont
clairement pas indépendantes mais on a

\[\text{Cov}(X,Y) = \text{Cov}(X,X^2) = \mathbb{E}(X^3) - \mathbb{E}(X)\mathbb{E}(X^2) = 0\]

\chapter{Identification de
densité}\label{II-identification-de-densituxe9-1}

Un problème important est le suivant. Soit \(X\) une variable aléatoire
réelle, admettant la densité \(f_X\). Soit \(g\) une fonction mesurable,
de sorte que \(Y = g(X)\) est aussi une variable aléatoire. Est-ce que
\(Y\) admet une densité, et si oui, comment la calculer\,?

On peut déjà remarquer que cette densité n'existe pas toujours. Si par
exemple \(g(x) = a\) pour tout \(x\), la loi de \(Y\) est la mesure de
Dirac en \(a\), qui n'a pas de densité par rapport à la mesure de
Lebesgue.

Pour résoudre ce problème, l'idée consiste à essayer de mettre
\(E(h(Y )) = E(h \circ g(X))\) sous la forme \(\int h(y)f_Y (y)dy\) pour
une fonction convenable \(f_Y\), et une classe de fonctions \(h\)
suffisamment grande. La fonction \(f_Y\) sera alors la densité cherchée.

La \protect\hyperref[II-esperanceg]{proposition qui assure que \(g(X)\)
est intégrable pour une variable aléatoire réelle
\(X\) (p.~\pageref*{II-esperanceg})} implique
\[ \mathbb{E}(h(Y)) = \mathbb{E}(h \circ g (X)) = \int_{\mathbb{R}} h \circ g(x) f_X(x) dx\]

et on fait le changement de variable \(y = g(x)\) dans cette intégrale.
Cela nécessite que \(g\) soit dérivable et bijective ``par morceaux'',
et il faut faire très attention aux domaines où \(g\) est croissante ou
décroissante. Puisque la fonction \(h\) est arbitraire, on appelle
couramment cette technique la \emph{méthode de la fonction muette}.
Cette approche résulte en fait de la proposition suivante~:

\addcontentsline{toc}{subsection}{Méthode de la fonction muette}
\begin{proposition}[{Méthode de la fonction
muette}]\phantomsection\label{II-muxe9thode-de-la-fonction-muette}

Si il existe une fonction \(f\) positive telle que pour toute fonction
mesurable \(h\) telle que \(h(x) f(x)\) soit intégrable,
\[\mathbb{E}(h(X)) = \int_\mathbb{R}h(x) f(x) dx\] alors la loi de \(X\)
admet la densité \(f\) par rapport à la mesure de Lebesgue.
\end{proposition}

\begin{preuve}
Supposons qu'un tel \(f\) existe. Soit
\(B \in \mathcal{B}(\mathbb{R})\), alors \(1_B\) est mesurable. En
particulier, si \(B = ]-\infty,x], ~x\in\mathbb{R}\), on a
\[\mathbb{P}(X\leq x) = \mathbb{P}(X \in B) = \mathbb{P}_X(B) = \mathbb{E}(1_B(X)) = \int_{- \infty}^x f(y)dy\]
Par ailleurs, prenant \(B = \mathbb{R}\), on a
\(1 = \mathbb{P}(X\in \mathbb{R}) = \int_\mathbb{R}f(y) dy\). Par
conséquent, \(\mathbb{P}_X\) admet la densité \(f\).\hfill$\blacksquare$
\end{preuve}

Nous donnons ici quelques exemples d'application de cette méthode~:

\begin{itemize}
\tightlist
\item
  Soit \(Y = aX + b\), où \(a\) et \(b\) sont des constantes. Si
  \(a=0\), alors \(Y = b\) et la loi de \(Y\) est la masse de Dirac en
  \(b\) (sans densité). Si \(a \neq 0\), on fait le changement de
  variable \(y = ax+b\), ce qui donne
  \[ \mathbb{E}(h(Y)) = \int_\mathbb{R}h(ax+b) f_X(x) dx = \int_\mathbb{R}h(y) f_X(\frac{y-b}{a})\frac{1}{|a|} dy \]
  Donc \[ f_Y(y) = f_X(\frac{y-b}{a})\frac{1}{|a|} \]
\item
  Soit \(Y =X^2\). La fonction \(g\) est décroissante sur
  \(\mathbb{R}_-\) et croissante sur \(\mathbb{R}_+\). Le changement de
  variable \(y = x^2\) donne alors \begin{align*}
       \mathbb{E}(h(Y)) &= \int_{-\infty}^0 h(x^2) f_X(x) dx + \int_0^{+\infty} h(x^2) f_X(x) dx\\
                  &= \int_0^{+\infty} h(y) f_X(-\sqrt{y})\frac{1}{2\sqrt{y}} dy + \int_0^{+\infty} h(y) f_X(\sqrt{y})\frac{1}{2\sqrt{y}}dy
       \end{align*} et on en déduit
  \[f_Y(y) = (f_X(-\sqrt{y})+f_X(\sqrt{y}))\frac{1}{2\sqrt{y}}1_{\left]0,+\infty\right[}\]
\end{itemize}

Dans le cas des vecteurs aléatoires, l'idée est la même. Soit
\(X = (X_1,\ldots,X_n)\), un vecteur aléatoire de densité \(f_X\) sur
\(\mathbb{R}^n\), \(g\) une fonction de \(\mathbb{R}^n\) dans
\(\mathbb{R}^m\) et \(Y = g(X)\). Plusieurs cas sont à considérer~:

\begin{enumerate}
\def\labelenumi{\arabic{enumi}.}
\tightlist
\item
  \(m > n\), le vecteur \(Y\) n'admet pas de densité.
\item
  \(m=n\), on utilise comme dans le cas unidimensionel le changement de
  variable \(y = g(x)\) dans
  \[ \mathbb{E}(h(Y)) = \mathbb{E}(h \circ g (X)) = \int_{\mathbb{R}^n} h \circ g(x) f_X(x) dx \]
  Supposons d'abord que \(g\) soit une bijection continûment
  différentiable de \(A\) dans \(B\), ouverts de \(\mathbb{R}^n\). Sous
  l'hypothèse que \(h \circ g(x) f_X(x)\) soit intégrable, le
  \href{Calcul\%20Intégral\%20III.pdf}{théorème de changement de
  variable} nous assure~:
  \[ \int_A h\circ g (x) f_X(x) dx = \int_B h(y) f_X \circ g^{-1}(y) \frac{1}{|\det J_g (g^{-1}(y))|} dy,\]
  où \(J_g\) désigne la matrice de Jacobi associée à la différentielle
  de \(g\). Dans le cas où \(f_X(x) = 0\) en dehors de \(A\), on obtient
  que \(Y\) admet la densité
  \[ f_Y(y) = 1_B(y)f_X \circ g^{-1}(y)\frac{1}{|\det J_g (g^{-1}(y))|}.\]
  Lorsque \(g\) est simplement continûment différentiable, il existe
  souvent une partition finie \((A_i)_{1\leq i \leq n}\) de l'ensemble
  \(\{x ; f(x) >0\}\), telle que \(g\) soit bijective sur chaque
  \(A_i\). On note alors \(B_i = g(A_i)\) l'image de \(A_i\) par \(g\).
  On découpe alors l'intégrale selon les \(A_i\), on applique la formule
  précédente à chaque morceau et on somme pour obtenir~:
  \[ f_Y(y) = \sum_{i=1}^n 1_{B_i}(y)f_X \circ g^{-1}(y) \frac{1}{|\det J_g (g^{-1}(y))|},\]
  où \(g^{-1}\) est bien définie sur chaque \(B_i\) comme image
  réciproque de la restriction de g à \(A_i\).
\item
  \(m < n\), on commence par ``compléter'' \(Y\), en essayant de
  construire une application \(g'\) de \(\mathbb{R}^n\) dans
  \(\mathbb{R}^n\) dont les \(m\) premières composantes coïncident avec
  les composantes de \(g\) et pour laquelle on peut appliquer l'une des
  deux formules précédentes. On obtient ainsi la densité \(f_Y'\) de
  \(Y' = g'(X)\) puis on obtient la densité de \(Y\) en calculant sa loi
  marginale~:
  \[f_Y(y_1,\ldots,y_m) = \int_{\mathbb{R}^{n-m}} f_{Y'}(y_1,\ldots,y_m,y_{m+1},\ldots y_n) dy_{m+1}\ldots dy_n.\]
\end{enumerate}

\addcontentsline{toc}{subsection}{\textbf{Coordonnées polaires}}
\begin{exemple}[{\textbf{Coordonnées
polaires}}]\phantomsection\label{II-polar}

Soit \(X = (U,V)\) un vecteur aléatoire de \(\mathbb{R}^2\), et
\(Y = (R,\Theta)\) ses coordonnées polaires. La transformation \(g\) est
un difféomorphisme de \(A = \mathbb{R}^2\setminus\{0\}\) dans
\(B = \left]0,+\infty\right[\times \left]0,2\pi\right]\), et son inverse
\(g^{-1}\) s'écrit~: \(u = r\cos \theta,\, v= r\sin \theta\). Le
Jacobien de \(g^{-1}\) au point \((r,\theta)\) vaut \(r\), donc le point
2. ci-dessus entraîne que
\[ f_Y(r,\theta) = r f_X(r\cos\theta,r\sin\theta)1_B(r,\theta)\] Par
exemple, si \(U\) et \(V\) sont indépendantes et de loi
\(\mathcal{N}(0,1)\), on a
\(f_X(u,v) = \frac{1}{2\pi}\exp\left(-\frac{u^2+v^2}{2}\right)\) et donc
\[ f_Y(r,\theta) = \frac{1}{2\pi} r\exp\left(-\frac{r^2}{2}\right)1_{\left]0,+\infty\right[ }(r)1_{\left]0,2\pi\right]}(\theta).\]
En particulier, on remarque que \(R\) et \(\Theta\) sont indépendantes
de densités respectives
\(r\exp\left(-\frac{r^2}{2}\right)1_{\left]0,+\infty\right[}(r)\) et
\(1_{\left]0,2\pi\right]}(\theta)\).
\end{exemple}

\runinhead{Exercice -- Lancer de fléchette
(fin) (\(\mathord{\bullet}\))}\phantomsection\addcontentsline{toc}{subsubsection}{Lancer
de fléchette (fin)}\label{II-exo-flechetteter}
Donner la densité du couple \((R,\Theta)\) des coordonnées polaires de
\(M\). Sont-elles indépendantes\,?\solref{II-answer-exo-flechetteter}

\addcontentsline{toc}{subsection}{\textbf{Somme de deux variables
aléatoires indépendantes}}
\begin{exemple}[{\textbf{Somme de deux variables aléatoires
indépendantes}}]\phantomsection\label{II-sum}

Soient \(U\) et \(V\) indépendantes et admettant les densités \(f_U\) et
\(f_V\), on cherche la densité de la somme \(Z = U+V\). On commence par
compléter \(Z\) en le couple \(T = (U,Z)\) (par exemple), correspondant
à la bijection \(g(u,v) = (u, u+v)\) sur \(\mathbb{R}^2\) dont le
jacobien est 1 et d'inverse \(g^{-1}(x,y) = (x,x-y)\). Appliquant le
point 3. ci-dessus, on obtient~:
\[f_Z(z) = \int_{\mathbb{R}}f_U(u)f_V(z-u) du = \int_{\mathbb{R}}f_U(z-v)f_V(v)dv.\]
La fonction \(f_Z\) est appelée \emph{le produit de convolution} des des
fonctions \(f_U\) et \(f_V\).
\end{exemple}

\runinhead{Exercice -- Loi
bêta (\(\mathord{\bullet}\mathord{\bullet}\mathord{\bullet}\))}\phantomsection\addcontentsline{toc}{subsubsection}{Loi
bêta}\label{II-exo-loibeta}
Soit \(X = (U,V)\) un vecteur aléatoire de \(\mathbb{R}^2\), avec \(U\)
et \(V\) indépendantes de lois \(\Gamma(\alpha,\theta)\) et
\(\Gamma(\beta,\theta)\). Identifier la densité de
\(Y = \frac{U}{U+V}\), puis de
\(Z = U+V\).\solref{II-answer-exo-loibeta}

\chapter{Modélisation statistique}\label{II-moduxe9lisation-statistique}

On a vu un certain nombre d'exemples de lois de probabilité à densité.
Celles-ci font intervenir des paramètres qui permettent d'ajuster la
forme de la densité aux données observées. Par exemple, la loi normale
fait intervenir deux paramètres~: l'espérance et la variance. Dans cette
section, nous formalisons l'approche statistique qui consiste à choisir
une loi de probabilité pour modéliser un phénomène aléatoire, puis à
estimer les paramètres de cette loi à partir d'observations.

La première hypothèse que nous faisons est que le phénomène étudié peut
être modélisé par une variable aléatoire \(X\) de densité \(f\)
inconnue. La première étape de la modélisation consiste à choisir une
famille de densités \((f_\theta)_{\theta \in \Theta}\), où \(\Theta\)
est un ensemble de paramètres. Par exemple, pour la loi normale, on peut
choisir \(\Theta = \mathbb{R}\times \mathbb{R}_+^*\) et
\(f_{(m,\sigma)}\) la densité de la loi \(\mathcal{N}(m,\sigma^2)\).

La seconde étape consiste à estimer le paramètre \(\theta\) de manière à
ce que la densité \(f_\theta\) soit la plus proche possible de la
densité \(f\) inconnue. On a donc besoin d'une notion de proximité entre
deux densités. Plusieurs choix sont possibles, mais le plus courant est
celui de la divergence de Kullback-Leibler.

\addcontentsline{toc}{subsection}{Divergence de Kullback-Leibler}
\begin{definition}[{Divergence de
Kullback-Leibler}]\phantomsection\label{II-defKL}

Soient \(f\) et \(g\) deux densités de probabilité sur \(\mathbb{R}\),
telles que \(g(x) > 0\) pour tout \(x\) tel que \(f(x) > 0\). La
\emph{divergence de Kullback-Leibler} de \(g\) par rapport à \(f\) est
définie par
\[ D_{KL}(f||g) = \int f(x) \log\left(\frac{f(x)}{g(x)}\right) dx.\]
\end{definition}

\begin{remarque}\phantomsection\label{II-remarque-1}
La divergence de Kullback-Leibler n'est pas une distance au sens
mathématique du terme. En effet, elle n'est pas symétrique et ne
satisfait pas l'inégalité triangulaire. En revanche, on a toujours
\(D_{KL}(f||g) \geq 0\), avec égalité si et seulement si \(f = g\)
presque partout.
\end{remarque}

\addcontentsline{toc}{subsection}{Positivité de la divergence de
Kullback-Leibler}
\begin{proposition}[{Positivité de la divergence de
Kullback-Leibler}]\phantomsection\label{II-corKL}

Soient \(f\) et \(g\) deux densités de probabilité sur \(\mathbb{R}\),
telles que \(g(x) > 0\) pour tout \(x\) tel que \(f(x) > 0\). Alors
\[ D_{KL}(f||g) \geq 0,\] avec égalité si et seulement si \(f = g\)
presque partout.
\end{proposition}

\begin{preuve}
La fonction \(\log\) est strictement concave, donc \(-\log\) est
strictement convexe. On peut donc appliquer
\protect\hyperref[II-jensen]{l'inégalité de
Jensen (p.~\pageref*{II-jensen})}~: \begin{align*}
D_{KL}(f||g) & = \int_\mathbb{R}f(x) \log\left(\frac{f(x)}{g(x)}\right) dx = - \int_\mathbb{R}f(x) \left(-\log\left(\frac{g(x)}{f(x)}\right)\right) dx \\
& = - \mathbb{E}\left(-\log\left(\frac{g(X)}{f(X)}\right)\right) \geq - \left(-\log\left(\mathbb{E}\left(\frac{g(X)}{f(X)}\right)\right)\right) = -\log(1) = 0.
\end{align*} On a égalité si et seulement si \(\frac{g(x)}{f(x)}\) est
presque sûrement constant, ce qui revient à dire que \(f = g\) presque
partout.\hfill$\blacksquare$
\end{preuve}

\runinhead{Exercice -- Divergence de Kullback-Leibler entre deux
gaussiennes}\phantomsection\addcontentsline{toc}{subsubsection}{Divergence
de Kullback-Leibler entre deux gaussiennes}\label{II-exoklgauss}
Soient \(f\) et \(g\) les densités de deux lois normales
\(\mathcal{N}(m_1,\sigma_1^2)\) et \(\mathcal{N}(m_2,\sigma_2^2)\).
Montrer que
\(D_{KL}(f||g) = \log \left(\frac{\sigma_2}{\sigma_1}\right) + \frac{\sigma_1^2 + (\mu_1 - \mu_2)^2}{2\sigma_2^2} - \frac{1}{2}\).\solref{II-answer-exoklgauss}

Revenons à notre problème d'estimation de paramètres. Soit \(X\) une
variable aléatoire de densité inconnue \(f\). On choisit une famille de
densités \((f_\theta)_{\theta \in \Theta}\), et on cherche à minimiser
la divergence de Kullback-Leibler entre \(f\) et \(f_\theta\)~:
\[ \hat{\theta} = \arg\min_{\theta \in \Theta} D_{KL}(f||f_\theta).\] On
peut réécrire la divergence de Kullback-Leibler sous la forme suivante~:
\begin{align*}
D_{KL}(f||f_\theta) & = \int f(x) \log\left(\frac{f(x)}{f_\theta(x)}\right) dx \\
& = \int f(x) \log(f(x)) dx - \int f(x) \log(f_\theta(x)) dx \\
& = -H(f) - \mathbb{E}_f(\log f_\theta(X)),
\end{align*} où \(H(f) = -\int f(x) \log f(x) dx\) est
l'\textbf{entropie} de la densité \(f\). Celle-ci s'interprète comme Ce
terme ne dépend pas de \(\theta\), donc minimiser
\(D_{KL}(f||f_\theta)\) revient à maximiser
\(\mathbb{E}_f(\log f_\theta(X))\). En statistique, \(f_\theta(X)\) est
appelée la \emph{vraisemblance} de l'observation \(X\) pour le paramètre
\(\theta\), et \(\log f_\theta(X)\) la \emph{log-vraisemblance}.
L'estimateur \(\hat{\theta}\) est ainsi appelé \emph{estimateur du
maximum de vraisemblance}.

En pratique, on dispose d'un échantillon \(X_1,\ldots,X_n\) de variables
aléatoires indépendantes et identiquement distribuées selon la loi de
densité \(f\). Le choix de la forme paramétrique est généralement guidé
par des considérations théoriques, portant sur la physique du phénomène
étudié, ou empiriques, en regardant par exemple l'histogramme des
données. On peut ensuite estimer \(\mathbb{E}_f(\log f_\theta(X))\) par
la moyenne empirique \[ \frac{1}{n} \sum_{i=1}^n \log f_\theta(X_i),\]
et l'estimateur du maximum de vraisemblance s'écrit
\[ \hat{\theta}_n = \arg\max_{\theta \in \Theta} \frac{1}{n} \sum_{i=1}^n \log f_\theta(X_i).\]
Cela revient à remplacer \(f\) par la mesure empirique de l'échantillon
\(\tilde{p} = \frac{1}{n} \sum_{i=1}^n \delta_{X_i}\) dans la définition
de la divergence de Kullback-Leibler.

\chapter{Exercices
complémentaires}\label{II-exercices-compluxe9mentaires}

\section{Durée de vie}\label{II-duruxe9e-de-vie}

La durée de vie, exprimée en années, d'un circuit électronique est une
variable aléatoire \(T\) dont la fonction de répartition \(F\) est
définie par~: \[F (t) = (1 - e^{-t^2 /2} )1_{t\geq 0}\]

\runinhead{Question
1}\phantomsection\addcontentsline{toc}{subsubsection}{Question
1}\label{II-dureevie1}
Donner la densité de probabilité \(f\) de \(T\) . Calculer
\(\mathbb{E}(T)\).\solref{II-answer-dureevie1}

\runinhead{Question
2}\phantomsection\addcontentsline{toc}{subsubsection}{Question
2}\label{II-dureevie2}
Sachant que le circuit a déjà fonctionné durant 1 an, quelle est la
probabilité qu'il continue à fonctionner encore durant au moins 2 ans\,?
La loi est-elle sans mémoire\,?

Un équipement électronique E est composé de 10 circuits identiques et
indépendants. Au circuit \(i (1 \leq i \leq 10)\) est associée la
variable aléatoire \(X_i\) , avec \(X_i = 1\) si la durée de vie du
circuit \(i\) est inférieure à un an et \(X_i = 0\)
sinon.\solref{II-answer-dureevie2}

\runinhead{Question
3}\phantomsection\addcontentsline{toc}{subsubsection}{Question
3}\label{II-dureevie3}
Quelle est la loi du nombre \(N\) de circuits de E dont la durée de vie
est inférieure à 1 an\,?\solref{II-answer-dureevie3}

\runinhead{Question
4}\phantomsection\addcontentsline{toc}{subsubsection}{Question
4}\label{II-dureevie4}
L'équipement E est dit en série si la défaillance de l'un de ses
circuits entraîne sa défaillance. Quelle est alors la probabilité qu'il
soit défaillant avant 1 an\,?\solref{II-answer-dureevie4}

\runinhead{Question
5}\phantomsection\addcontentsline{toc}{subsubsection}{Question
5}\label{II-dureevie5}
L'équipement E est dit en parallèle si sa défaillance ne peut se
produire que si tous ses circuits sont défaillants. Quelle est alors la
probabilité qu'il soit défaillant avant 1
an\,?\solref{II-answer-dureevie5}

\section{Crues centennales de la
Seine}\label{II-crues-centennales-de-la-seine}

On suppose que la hauteur d'eau maximale au cours de l'année \(n\) est
décrite par une variable aléatoire \(X_n\) de densité \(f(x)\) et de
fonction de répartition \(F(x)\), identique pour toutes les \(X_n\), que
l'on suppose également indépendantes.

\runinhead{Question
1}\phantomsection\addcontentsline{toc}{subsubsection}{Question
1}\label{II-crue1}
Calculer la fonction de répartition du maximum sur \(n\) années
\(Y = \max(X_1,\ldots,X_n)\).\solref{II-answer-crue1}

\runinhead{Question
2}\phantomsection\addcontentsline{toc}{subsubsection}{Question
2}\label{II-crue2}
En déduire la densité de \(Y\).\solref{II-answer-crue2}

\runinhead{Question
3}\phantomsection\addcontentsline{toc}{subsubsection}{Question
3}\label{II-crue3}
Expliciter la densité de \(Y\) lorsque les \(X_i\) sont de loi uniforme
sur \([0,1]\).\solref{II-answer-crue3}

\runinhead{Question
4}\phantomsection\addcontentsline{toc}{subsubsection}{Question
4}\label{II-crue4}
Selon mêmes hypothèses, calculer la hauteur minimale des quais pour que
la probabilité de voir la Seine déborder sur une période de \(n\) années
soit inférieure à 1/1000.\solref{II-answer-crue4}

\runinhead{Question
5}\phantomsection\addcontentsline{toc}{subsubsection}{Question
5}\label{II-crue5}
Toujours selon les mêmes hypothèses, calculer la médiane de \(Y\). La
comparer avec celle des \(X_i\). Commenter.\solref{II-answer-crue5}

\section{Indépendance et vecteurs
gaussiens}\label{II-induxe9pendance-et-vecteurs-gaussiens}

Soit \(Z := (X,Y)\) un vecteur gaussien (i.e.~qui suit une loi Normale
bi-variée) d'espérance \(m\) et de matrice de covariance définie
positive \(C\). Notons \(f_Z\) sa densité.

\runinhead{Question
1}\phantomsection\addcontentsline{toc}{subsubsection}{Question
1}\label{II-covindepgauss-joint}
Rappeler la forme de \(f_Z\) et en donner une expression non matricielle
faisant apparaître le coefficient de corrélation noté
\(\rho\).\solref{II-answer-covindepgauss-joint}

\runinhead{Question
2}\phantomsection\addcontentsline{toc}{subsubsection}{Question
2}\label{II-covindepgauss-marg}
Expliciter les lois marginales de \(X\) et de
\(Y\).\solref{II-answer-covindepgauss-marg}

\runinhead{Question
3}\phantomsection\addcontentsline{toc}{subsubsection}{Question
3}\label{II-covindepgauss-indep}
Montrer que \(X\) et \(Y\) sont indépendantes ssi
\(\text{Cov}\left(X,Y\right) = 0\).\solref{II-answer-covindepgauss-indep}

\section{Symétrie de la
gaussienne}\label{II-symuxe9trie-de-la-gaussienne}

Soit \(X\) une variable aléatoire réelle de loi Normale centrée réduite,
dont la densité est notée \(f\) et la fonction de répartition \(F\).
Pour tout \(c>0\) on pose
\[X_c := \left|\begin{array}{ll} X & \text{si } |X| > c,\\ -X & \text{si } |X| \leq c.\end{array}\right.\]

\runinhead{Question
1}\phantomsection\addcontentsline{toc}{subsubsection}{Question
1}\label{II-gausssym-loi}
Déterminer la loi de \(X_c\).\solref{II-answer-gausssym-loi}

\runinhead{Question
2}\phantomsection\addcontentsline{toc}{subsubsection}{Question
2}\label{II-gausssym-cov}
Calculer la covariance de \(X\) et de \(X_c\) et montrer qu'il existe
une valeur \(c_0\) de \(c\) qui l'annule.\solref{II-answer-gausssym-cov}

\runinhead{Question
3}\phantomsection\addcontentsline{toc}{subsubsection}{Question
3}\label{II-gausssym-indep}
Les deux variables \(X\) et \(X_{c_0}\) sont-elles
indépendantes\,?\solref{II-answer-gausssym-indep}

\runinhead{Question
4}\phantomsection\addcontentsline{toc}{subsubsection}{Question
4}\label{II-gausssym-vecgauss}
Le vecteur aléatoire \((X,X_{c_0})\) est-il
gaussien\,?\solref{II-answer-gausssym-vecgauss}

\section{Loi du \(\chi^2\)}\label{II-loi-du-chi2}

On considère une variable aléatoire \(X\) gaussienne centrée réduite. On
note \(f\) sa densité et \(F\) sa fonction de répartition.

\runinhead{A 1 degré de
liberté}\phantomsection\addcontentsline{toc}{subsubsection}{A 1 degré de
liberté}\label{II-khi2-1dl}

\begin{enumerate}
\def\labelenumi{\arabic{enumi}.}
\tightlist
\item
  Calculer la densité de la variable aléatoire \(Y := X^2\). En donner
  une expression faisant apparaître la fonction gamma (on rappelle que
  \(\Gamma\left(\frac{1}{2}\right)=\sqrt{\pi}\)).
\end{enumerate}

On dit que \(Y\) suit une \emph{loi du \(\chi^2\) à \(1\) degré de
liberté}, et on note \(Y \sim \chi^2\).\solref{II-answer-khi2-1dl}

\runinhead{A n degrés de
liberté}\phantomsection\addcontentsline{toc}{subsubsection}{A n degrés
de liberté}\label{II-khi2-ndl}
Soient maintenant \(n\in\mathbb{N}^\ast\) et \(X_1,\dots,X_n\) des
copies indépendantes de \(X\). Pour tout \(i \in \{1,\dots,n\}\) on pose
\(Y_i := X_i^2\).

\begin{enumerate}
\def\labelenumi{\arabic{enumi}.}
\setcounter{enumi}{1}
\tightlist
\item
  Montrer que la variable aléatoire \(Y := \sum_{i = 1}^n Y_i\) admet
  pour densité
  \[ f_Y : x\in\mathbb{R}\mapsto \left|\begin{array}{ll} \dfrac{x^{\frac{n}{2} - 1}}{2^{\frac{n}{2}}\,\Gamma\left(\frac{n}{2}\right)}\,\exp\left\{ -\dfrac{x}{2} \right\} & \text{si } x> 0,\\ 0 & \text{si } x \leq 0. \end{array}\right.\]
  On pourra procéder par récurrence sur \(n\).
\end{enumerate}

On dit que \(Y\) suit une \emph{loi du \(\chi^2\) à \(n\) degrés de
liberté} et on note \(Y \sim \chi^2_n\).\solref{II-answer-khi2-ndl}

\section{Combinaisons linéaires de variables Gaussiennes
indépendantes}\label{II-combinaisons-linuxe9aires-de-variables-gaussiennes-induxe9pendantes}

Soit \(X\) une variable aléatoire de loi Normale centrée réduite, dont
on note \(f\) la densité et \(F\) la fonction de répartition.

\runinhead{Préliminaires}\phantomsection\addcontentsline{toc}{subsubsection}{Préliminaires}\label{II-CLIGauss-pre}

\begin{enumerate}
\def\labelenumi{\arabic{enumi}.}
\tightlist
\item
  Rappeler la loi de la variable aléatoire \(s\,X + m\), où
  \(s, m \in \mathbb{R}\).
\end{enumerate}

\solref{II-answer-CLIGauss-pre}

\runinhead{Combinaisons
linéaires}\phantomsection\addcontentsline{toc}{subsubsection}{Combinaisons
linéaires}\label{II-CLIGauss-cl}
Soient maintenant \(n\in\mathbb{N}^\ast\) variables aléatoires
\(X_1,\dots,X_n\) indépendantes et de même loi que \(X\) (on dit que ce
sont des \emph{copies indépendantes} de \(X\)). Pour tout vecteur
\(a = (a_1,\dots,a_n) \in (\mathbb{R}^{*})^n\) on pose
\(S_n^a := \sum_{i = 1}^n a_i\,X_i\).

\begin{enumerate}
\def\labelenumi{\arabic{enumi}.}
\setcounter{enumi}{1}
\item
  Montrer que pour \(S_n^a\) suit une loi Normale d'espérance nulle et
  de variance \(\sum_{i = 1}^n a_i^2\) (on pourra raisonner par
  récurrence sur \(n\)).
\item
  Soient \(a,b\in(\mathbb{R}^{*})^n\). Sous quelle condition a-t-on
  \(\text{Cov}\left(S_n^a, S_n^b\right) = 0\)\,?
\end{enumerate}

\solref{II-answer-CLIGauss-cl}

\section{Loi de Maxwell}\label{II-loi-de-maxwell}

On désire déterminer la distribution des vitesses des molécules d'un gaz
monoatomique parfait à l'équilibre (loi de Maxwell (1859)).

On représente la vitesse d'une molécule d'un gaz monoatomique parfait à
l'équilibre dans un repère orthonormal par un vecteur aléatoire
\(V = (V_1 , V_2 , V_3 )\). Le choix du repère étant arbitraire, il est
naturel de supposer que la loi de \(V\) est invariante par rotation
(autour de l'origine) et que les composantes de \(V\) sont
indépendantes.

\subsection{Partie 1}\label{II-partie-1}

Soit \((X, Y, Z)\) un vecteur aléatoire à valeurs dans \(\mathbb{R}^3\)
de densité \(f_{(X, Y, Z)}\). Ce vecteur aléatoire est supposé invariant
par rotation~: il existe une fonction
\(\phi: \left[0, +\infty\right[ \to \left]0, +\infty\right[\)
(strictement positive) telle que \[
f_{(X, Y, Z)}(x,y,z) = \phi(x^2+y^2+z^2) \;  \mbox{ pour tout $(x, y, z) \in \mathbb{R}^3$.}
\]

\runinhead{Question
1}\phantomsection\addcontentsline{toc}{subsubsection}{Question
1}\label{II-Maxwell-1}
On suppose de plus les variables aléatoires \(X\), \(Y\) et \(Z\)
indépendantes. Montrer qu'il existe une fonction
\(f: \mathbb{R}\to \left[0, +\infty\right[\) qui soit une densité telle
que \[
f_{(X, Y, Z)}(x, y, z) = f(x) f(y) f(z) \; \mbox{ pour tout $(x, y, z) \in \mathbb{R}^3$.} 
\] Montrer finalement que \(f(x) > 0\) pour tout
\(x \in \mathbb{R}\).\solref{II-answer-Maxwell-1}

\runinhead{Question
2}\phantomsection\addcontentsline{toc}{subsubsection}{Question
2}\label{II-Maxwell-2}
On suppose de plus que les densités marginales de \(X, Y\) et \(Z\)
ainsi que la fonction \(\phi\) sont continûment différentiables
(c'est-à-dire de classe \(C^1\)). Montrer que la densité de chacune des
composantes s'écrit sous la forme \(f (x) = a e^{cx^2/2}\), avec \(a\)
et \(c\) deux constantes réelles.\solref{II-answer-Maxwell-2}

\runinhead{Question
3}\phantomsection\addcontentsline{toc}{subsubsection}{Question
3}\label{II-Maxwell-3}
En déduire que le vecteur \((X, Y, Z)\) suit une loi gaussienne
d'espérance l'origine dont on précisera la matrice de covariance en
fonction de la constante
\(\sigma = 1/\sqrt{|c|}\).\solref{II-answer-Maxwell-3}

\subsection{Partie 2}\label{II-partie-2}

On suppose que le vecteur aléatoire \((X, Y, Z) = V\) vérifie les
hypothèses des questions précédentes.

\runinhead{Question
4}\phantomsection\addcontentsline{toc}{subsubsection}{Question
4}\label{II-Maxwell-4}
Calculer l'énergie cinétique moyenne d'un atome du gaz, c'est-à-dire
l'espérance \[
E_c := \mathbb{E}\left(\frac{1}{2}m\|V\|^2\right)
\] où \(m\) est la masse d'un atome du gaz. L'énergie cinétique moyenne
d'un atome du gaz de masse \(m\) étant égale à \(\frac{3}{2} kT\) où
\(k\) est la constante de Boltzmann et \(T\) la température du gaz, en
déduire la valeur de \(\sigma^2\) en fonction de \(k\), \(T\) et
\(m\).\solref{II-answer-Maxwell-4}

\runinhead{Question
5}\phantomsection\addcontentsline{toc}{subsubsection}{Question
5}\label{II-Maxwell-5}
On rappelle que si \(X\) et \(Y\) sont deux variables aléatoires
indépendantes de loi respective \(\Gamma (a,\lambda)\) et
\(\Gamma (b,\lambda)\), alors la loi de \(X + Y\) est la loi
\(\Gamma (a+b,\lambda)\). On rappelle également que la densité \(g\) de
la loi \(\Gamma (a,\lambda)\) vérifie \[
g(x) = \frac{1}{\Gamma(a)}\lambda^a x^{a -1} e^{-\lambda x} 1_{]0,+\infty[}(x)
\; \mbox{ pour tout $x\in\mathbb{R}$},
\] où \(\Gamma\) est la fonction définie par
\[\Gamma(a) = \int_0^{+\infty} x^{a-1}e^{-x} dx \text{ pour } a \in \left] 0, + \infty \right[.\]
On remarquera que \(\Gamma(1/2) = \sqrt{\pi}\) et que
\(\Gamma(x+1) = x \Gamma(x).\)

Calculer la loi de \(V_1^2\). En déduire la loi de \(\|V \|^2\) puis la
densité de \(\|V \| = \sqrt{V_1^2 + V_2^2 + V_3^2}\). La probabilité
associée est appelée loi de Maxwell.

\begin{center}\rule{0.5\linewidth}{0.5pt}\end{center}

\solref{II-answer-Maxwell-5}
